% Generated draft: requires mathematical and source review.
\documentclass[11pt]{article}
\usepackage[margin=25mm]{geometry}
\usepackage{fontspec}
\setmainfont{Latin Modern Roman}
\usepackage{amsmath,amssymb,mathtools,unicode-math}
\setmathfont{Latin Modern Math}
\usepackage{enumitem,needspace,xurl,hyperref}
\setlength{\parindent}{0pt}
\setlength{\parskip}{5pt}
\setlength{\emergencystretch}{3em}
\newcommand{\N}{\mathbb{N}}
\newcommand{\Z}{\mathbb{Z}}
\newcommand{\Q}{\mathbb{Q}}
\newcommand{\R}{\mathbb{R}}
\newcommand{\C}{\mathbb{C}}
\newcommand{\F}{\mathbb{F}}
\newcommand{\A}{\mathbb{A}}
\newcommand{\PP}{\mathbb{P}}
\newcommand{\E}{\mathbb{E}}
\newcommand{\eps}{\varepsilon}
\newcommand{\dd}{\,\mathrm{d}}
\begin{document}
\section*{UnsolvedMath ID: 8800073}
\textbf{Source problem code:} catalogue entry
\subsection{Definitions and mathematical statement}
Fix block length \(m\) and block number \(\ell\), so \(N=m\ell\). A standard quasi-cyclic symmetry simultaneously rotates the \(m\) coordinates in each of \(\ell\) blocks. A linear code \(C\le\F_q^N\) is quasi-cyclic when invariant under that permutation; a full-rank lattice \(L\le\Z^N\) is quasi-cyclic under the analogous integer-coordinate symmetry. Equivalently one obtains module structures over suitable quotients by \(x^m-1\), with the ring and block conventions fixed. This is not automatically the same family as ideals in a chosen cyclotomic number field.

The programme is to construct/analyze an appropriate public-key primitive using these structures and prove its chosen security property by reduction from a worst-case problem on the intended quasi-cyclic class. Candidate underlying problems could be approximate lattice problems or bounded-weight syndrome decoding
\[
He=s,\qquad\operatorname{wt}(e)\le w,
\]
for a specified quasi-cyclic code representation. The source does not select one.

The required reduction should transform an efficient adversary with nonnegligible advantage against keys/challenges drawn from the scheme's actual distribution into an efficient solver for every input of the chosen worst-case structured problem, with stated approximation and success losses. Thus worst-case hardness of that problem would imply security in the specified experiment. Merely assuming that random structured instances are difficult, or citing unrelated worst-case NP-hardness, is not this conclusion. Nor is the question an assertion that every conceivable quasi-cyclic public-key construction is secure.

Parameters \(m,\ell,q\), underlying hardness problem and approximation factor, key distribution, security notion, and classical/quantum reduction model are absent. The supplied research summary reports no verified general clean reduction of the requested kind and mentions BIKE/HQC and Ring-LWE as related examples. Ring-LWE results for specified ideal/module settings cannot be transferred to arbitrary quasi-cyclic codes by terminology alone. An explicit average-case assumption can still support a rigorous conditional proof; it is not synonymous with an unformalized heuristic.
\subsection{Short English statement}
Give a quasi-cyclic lattice- or code-based public-key construction with an explicit security reduction from a specified worst-case structured problem. Fix the structure, approximation parameters, distributions, and security model rather than relying only on average-case difficulty.
\subsection{Sources}
\begin{itemize}
\item Local UnsolvedMath registry, permanent ID \(8800073\), exact source code \texttt{catalogue entry}, ``Worst-case security of quasi-cyclic cryptosystems.'' Proposer Daniele Micciancio.
\item IPAM \emph{Workshop on Number Theory and Cryptography --- Open Problems} (2006), attachment 6240, recorded PDF page 39; \url{https://www.ipam.ucla.edu/programs/workshops/workshop-i-number-theory-and-cryptography-open-problems/?tab=schedule}. Recorded access 2026-07-29; pdf-and-legacy-ppt visual-audit extraction.
\item Supplied OPEN-TRIAGE summary mentions BIKE, HQC, and Ring-LWE but provides no full citations, exact reduction theorem, or matching quasi-cyclic hardness/security parameters.
\end{itemize}
\end{document}
